\begin{figure}[tbp]
\centering
\begin{tikzpicture}[
  >=Latex,
  leg/.style={circle,draw,minimum size=7mm,inner sep=0pt},
  every node/.style={font=\small}
]
  \node (sum) at (-4.7,0.7) {$\displaystyle T=\sum_j\lambda_j$};
  \node[leg] (a) at (-3.2,1.45) {$e_{1j}$};
  \node[leg] (b) at (-2.15,0.7) {$e_{2j}$};
  \node[leg] (c) at (-3.2,-0.05) {$e_{3j}$};
  \draw[very thick] (a) -- (b) -- (c) -- cycle;
  \node[font=\scriptsize,fill=white,inner sep=1pt] at (-2.95,0.7) {$j$};

  \draw[dashed] (-0.55,-1.45) -- (-0.55,2.0);

  \node[align=center] at (1.2,1.55) {generic contraction\\$x=\sum_jx_je_{1j}^*$};
  \draw[step=0.32,black!35] (0.18,-0.1) grid (1.78,1.18);
  \foreach \x/\y in {0.34/0.06,0.66/0.38,0.98/0.70,1.30/1.02}
    \fill (\x,\y) rectangle ++(0.32,0.32);
  \node[align=center] at (1.0,-0.58)
    {$\rank(T_x)=\#\{j:x_j\ne0\}$};

  \node[align=center] at (4.25,1.55) {axis contraction\\$x=e_{1k}^*$};
  \draw[step=0.32,black!35] (3.45,-0.1) grid (5.05,1.18);
  \fill (4.09,0.54) rectangle ++(0.32,0.32);
  \node at (4.25,-0.58) {$\rank(T_x)=1$};

  \draw[decorate,decoration={brace,amplitude=5pt},thick]
    (0.15,-1.05) -- (5.1,-1.05)
    node[midway,below=6pt,align=center]
      {rank-one contraction locus $=$ dual coordinate axes\\
       product equivalence must permute the axes};
\end{tikzpicture}
\caption{A diagonal tensor with all coefficients nonzero exposes its own
axes.  Contracting one leg against \(x\) and flattening the remainder gives
rank equal to the number of nonzero coordinates of \(x\), so the rank-one
contraction locus is exactly the set of dual coordinate axes.  For an
\((m+1)\)-party AME marginal in local Hilbert--Schmidt spaces, the recovered
axes are the local Weyl axes, and a product conjugation must permute them.}
\label{fig:axis-recovery}
\end{figure}
